\section{Step 2: Question-Level Behavioral Harmonization and Descriptive Analysis}
\label{sec:step2_behavior}
% Analyst note added 2026-03-20: PID006 and PID037 are globally excluded from downstream participant-level analyses. PID006 was excluded after Step 3 documented session 2025-11-12_002 as an implausibly short 10.420224 second recording with no paired PsychoPy log. PID037 was excluded because a behavioral log exists but no fNIRS session folder was found. Step 2 outputs are regenerated without these participants, and future steps should honor materials/analysis_participant_exclusions.json.

\subsection{Objective}
The objective of Step~2 is to build a clean, question-level behavioral dataset that links:
\begin{itemize}
    \item participant responses from the PsychoPy task logs;
    \item question metadata from \texttt{filtered\_questions.json};
    \item derived behavioral measures such as response time and correctness;
    \item descriptive comparisons between participant performance and the original ENEM item correctness.
\end{itemize}

This step remains intentionally simple.
It focuses only on the behavioral side of the task and does not yet modify fNIRS files or analyze the fNIRS signals.
Its purpose is to produce a well-organized and fully documented behavioral dataset that can later be aligned with the fNIRS recordings.

\subsection{Main Goal}
In practical terms, Step~2 should answer the following basic questions:
\begin{itemize}
    \item Which questions were shown to each participant?
    \item Which answers did participants choose?
    \item Were those answers correct or incorrect?
    \item How long did participants take on each question?
    \item How do participant-level and item-level performance patterns differ between abstract and concrete questions?
    \item How does the observed participant accuracy compare with the original ENEM correctness for the same items?
    \item How do simple text properties of the items, such as word count and sentence length, relate to the observed behavioral patterns?
\end{itemize}

\subsection{Raw Data Sources}
Step~2 relies on two main raw sources.

\subsubsection*{1. PsychoPy task log}
For each participant, the file
\begin{verbatim}
enem_blocks_PIDxx_YYYYMMDD_HHMMSS.csv
\end{verbatim}
contains event-level task information, including at least the following columns:
\begin{verbatim}
t_abs
phase
block
trial_idx_in_block
question_number
question_year
question_type
question_field
marker_name
marker_code
rt_from_phase
choice
correct
button_click_time
option_view_time
note
\end{verbatim}

The relevant phases include:
\begin{itemize}
    \item \texttt{q\_text\_on}
    \item \texttt{q\_stem\_on}
    \item \texttt{q\_options\_on}
    \item \texttt{answer}
    \item \texttt{iti}
    \item \texttt{block\_start}
    \item \texttt{block\_end}
\end{itemize}

\subsubsection*{2. Question metadata}
The file
\begin{verbatim}
filtered_questions.json
\end{verbatim}
contains item-level metadata, including variables such as:
\begin{verbatim}
type
correctness
answer
total_word_count
sentence_count
sentence_length
\end{verbatim}

Depending on the version of the file, it may also contain question text, field, year, options, and other metadata.

\subsection{Interpretation of the Raw Structure}
The PsychoPy task presents each item in stages:
\begin{enumerate}
    \item text onset;
    \item stem onset;
    \item options onset;
    \item answer.
\end{enumerate}

This is useful because Step~2 can reconstruct not only a single response time, but also a small set of simple timing components:
\begin{itemize}
    \item time from text onset to stem onset;
    \item time from stem onset to options onset;
    \item time from options onset to answer;
    \item total time from first presentation to answer.
\end{itemize}

This is still a basic descriptive approach, but it is more informative than keeping only one undifferentiated response-time measure.

\subsection{Scope}
This step includes:
\begin{itemize}
    \item extraction of question-level behavioral data from the PsychoPy logs;
    \item harmonization of question identifiers;
    \item creation of a clean participant-question table;
    \item derivation of timing variables;
    \item derivation of participant correctness by comparing selected answers with the correct option;
    \item joining participant responses with question metadata from \texttt{filtered\_questions.json};
    \item descriptive summaries by question, question type, block, field, year, and participant;
    \item descriptive comparison between participant sample accuracy and ENEM correctness;
    \item export of CSV outputs, plots, and modular \LaTeX{} files.
\end{itemize}

This step excludes:
\begin{itemize}
    \item trigger repair in the fNIRS files;
    \item fNIRS preprocessing;
    \item fNIRS statistical analysis;
    \item advanced inferential modeling;
    \item mixed models, hierarchical models, or psychometric models.
\end{itemize}

\subsection{Unit of Analysis}
Step~2 should explicitly define three units of analysis.

\subsubsection*{Participant-question level}
This is the main unit.
Each row corresponds to one participant responding to one question.

\subsubsection*{Question level}
Each row corresponds to one unique question, aggregated across all participants who saw that question.

\subsubsection*{Participant level}
Each row corresponds to one participant, summarizing their overall task behavior and their behavior separately for abstract and concrete questions.

\subsection{Core Principle}
Because different participants may have seen different questions, and because block timing may cause some questions not to be reached, every summary in Step~2 must always report its denominator.
In other words, all descriptive results must make clear:
\begin{itemize}
    \item how many participants were exposed to the item;
    \item how many answered it;
    \item how many valid response times are available;
    \item how many valid correctness labels are available.
\end{itemize}

\subsection{Recommended Question Identifier}
A stable question identifier should be defined before any merge.
A practical recommendation is:
\begin{verbatim}
question_id = question_year + "_" + question_field + "_" + question_number
\end{verbatim}

If \texttt{filtered\_questions.json} already contains a unique item ID, that identifier should be preferred.
The same identifier must be used in all Step~2 outputs.

\subsection{Derived Variables}
The clean trial table should contain, at minimum, the following groups of variables.

\subsubsection*{Participant and session variables}
\begin{itemize}
    \item \texttt{participant\_id}
    \item \texttt{timestamp}
    \item \texttt{source\_log}
\end{itemize}

\subsubsection*{Task-position variables}
\begin{itemize}
    \item \texttt{block}
    \item \texttt{block\_type} (Abstract or Concrete)
    \item \texttt{block\_index\_within\_type} (1 to 5)
    \item \texttt{global\_block\_order} (1 to 10)
    \item \texttt{trial\_idx\_in\_block} (1 to 3)
    \item \texttt{trial\_idx\_global}
    \item \texttt{order\_condition} (AbstractFirst or ConcreteFirst)
\end{itemize}

\subsubsection*{Question metadata variables}
\begin{itemize}
    \item \texttt{question\_id}
    \item \texttt{question\_number}
    \item \texttt{question\_year}
    \item \texttt{question\_field}
    \item \texttt{question\_type}
    \item \texttt{correct\_answer}
    \item \texttt{enem\_correctness}
    \item \texttt{total\_word\_count}
    \item \texttt{sentence\_count}
    \item \texttt{sentence\_length}
\end{itemize}

\subsubsection*{Behavioral response variables}
\begin{itemize}
    \item \texttt{chosen\_answer}
    \item \texttt{participant\_correct}
    \item \texttt{response\_status}
    \item \texttt{answer\_available}
\end{itemize}

\subsubsection*{Timing variables}
\begin{itemize}
    \item \texttt{t\_text\_on}
    \item \texttt{t\_stem\_on}
    \item \texttt{t\_options\_on}
    \item \texttt{t\_answer}
    \item \texttt{rt\_text\_to\_stem}
    \item \texttt{rt\_stem\_to\_options}
    \item \texttt{rt\_options\_to\_answer}
    \item \texttt{rt\_text\_to\_answer}
    \item \texttt{rt\_stem\_to\_answer}
\end{itemize}

\subsubsection*{Normalized timing variables}
Because questions differ in length, it is useful to create very simple normalized timing measures:
\begin{itemize}
    \item \texttt{sec\_per\_word} = \texttt{rt\_text\_to\_answer / total\_word\_count}
    \item \texttt{sec\_per\_sentence} = \texttt{rt\_text\_to\_answer / sentence\_count}
\end{itemize}

These variables remain descriptive only.
They are included because raw time alone may be misleading when items differ substantially in length.

\subsection{Expected Outputs}

\subsubsection*{Structured tables}
\begin{verbatim}
01_question_file_manifest.csv
02_question_metadata_raw.csv
03_question_metadata_clean.csv
04_behavior_events_raw.csv
05_behavior_trial_raw.csv
06_behavior_trial_clean.csv
07_question_summary.csv
08_question_answer_distribution.csv
09_type_summary.csv
10_block_summary.csv
11_field_summary.csv
12_year_summary.csv
13_participant_behavior_summary.csv
14_exam_vs_sample_comparison.csv
15_behavior_missingness.csv
\end{verbatim}

\subsubsection*{Plots}
\begin{verbatim}
figures/step2/accuracy_by_type.png
figures/step2/rt_total_by_type.png
figures/step2/enem_vs_sample_accuracy.png
figures/step2/word_count_vs_rt.png
figures/step2/word_count_vs_accuracy.png
figures/step2/question_accuracy_comparison.png
figures/step2/question_rt_distribution.png
\end{verbatim}

\subsubsection*{\LaTeX{} outputs}
\begin{verbatim}
reports/step2/questions/question_<question_id>.tex
reports/step2/groups/group_abstract_vs_concrete.tex
reports/step2/groups/group_blocks.tex
reports/step2/groups/group_fields.tex
reports/step2/step2_behavior_report.tex
\end{verbatim}

\subsection{Detailed Workflow}

\subsubsection{Step~2.1: Build the file manifest}
Create a manifest that lists, for each participant:
\begin{itemize}
    \item participant folder;
    \item PsychoPy log file name;
    \item timestamp extracted from file name;
    \item whether the log is readable;
    \item whether a paired demographic file exists;
    \item notes on any irregularity.
\end{itemize}

The purpose of this file is to ensure that all later behavioral analyses are traceable back to the raw source file.

\subsubsection{Step~2.2: Parse and flatten \texttt{filtered\_questions.json}}
Convert \texttt{filtered\_questions.json} into a tabular question-metadata file.
For each question, extract at least:
\begin{itemize}
    \item question ID;
    \item year;
    \item field;
    \item type;
    \item correct answer;
    \item ENEM correctness;
    \item total word count;
    \item sentence count;
    \item sentence length.
\end{itemize}

Save this as:
\begin{verbatim}
02_question_metadata_raw.csv
\end{verbatim}

\subsubsection{Step~2.3: Clean and validate question metadata}
Create a cleaned question metadata table.
At this stage:
\begin{itemize}
    \item standardize field names;
    \item standardize type labels to exactly \texttt{Abstract} and \texttt{Concrete} or a fixed lower-case equivalent;
    \item validate that ENEM correctness is numeric and lies between 0 and 1;
    \item validate that correct answer is one of \texttt{A}, \texttt{B}, \texttt{C}, \texttt{D}, or \texttt{E};
    \item validate that text-complexity variables are numeric and positive when present.
\end{itemize}

Save the cleaned table as:
\begin{verbatim}
03_question_metadata_clean.csv
\end{verbatim}

\subsubsection{Step~2.4: Extract event-level task data from the PsychoPy logs}
Read all \texttt{enem\_blocks} files and concatenate the relevant task rows into one raw event table.
Keep only the fields needed for behavioral reconstruction:
\begin{itemize}
    \item participant ID;
    \item timestamp;
    \item phase;
    \item block;
    \item trial index in block;
    \item question identifiers;
    \item \texttt{t\_abs};
    \item \texttt{choice};
    \item \texttt{button\_click\_time};
    \item \texttt{option\_view\_time};
    \item marker fields;
    \item notes.
\end{itemize}

Save this as:
\begin{verbatim}
04_behavior_events_raw.csv
\end{verbatim}

\subsubsection{Step~2.5: Reconstruct one row per participant-question trial}
Group the event table by participant, block, trial index, and question identifier, and reconstruct one trial row for each presented question.

For each trial, identify:
\begin{itemize}
    \item the text onset event;
    \item the stem onset event;
    \item the options onset event;
    \item the answer event, if present.
\end{itemize}

This yields a raw trial table with one row per participant-question attempt.
Save it as:
\begin{verbatim}
05_behavior_trial_raw.csv
\end{verbatim}

\subsubsection{Step~2.6: Derive timing variables}
For each reconstructed trial, compute:
\begin{itemize}
    \item \texttt{rt\_text\_to\_stem = t\_stem\_on - t\_text\_on}
    \item \texttt{rt\_stem\_to\_options = t\_options\_on - t\_stem\_on}
    \item \texttt{rt\_options\_to\_answer = t\_answer - t\_options\_on}
    \item \texttt{rt\_text\_to\_answer = t\_answer - t\_text\_on}
    \item \texttt{rt\_stem\_to\_answer = t\_answer - t\_stem\_on}
\end{itemize}

Also compute:
\begin{itemize}
    \item \texttt{sec\_per\_word}
    \item \texttt{sec\_per\_sentence}
\end{itemize}

If any phase is missing, the derived timing variable should be set to missing and flagged explicitly.

\subsubsection{Step~2.7: Derive answer and correctness variables}
The PsychoPy log stores the participant's selected option in the \texttt{choice} column of the \texttt{answer} row.
Use this selected option and compare it with the correct option from \texttt{filtered\_questions.json}.

Create:
\begin{itemize}
    \item \texttt{chosen\_answer}
    \item \texttt{correct\_answer}
    \item \texttt{participant\_correct}
\end{itemize}

Recommended coding:
\begin{itemize}
    \item \texttt{participant\_correct = 1} if chosen answer matches correct answer;
    \item \texttt{participant\_correct = 0} if chosen answer does not match correct answer;
    \item missing if the answer event is absent or unusable.
\end{itemize}

\subsubsection{Step~2.8: Define response-status categories}
Use simple explicit status labels:
\begin{itemize}
    \item \texttt{ANSWERED\_VALID}
    \item \texttt{ANSWERED\_INVALID}
    \item \texttt{MISSING\_ANSWER}
    \item \texttt{PARTIAL\_TRIAL}
    \item \texttt{UNMATCHED\_QUESTION}
\end{itemize}

These categories should be used throughout Step~2 to keep the data traceable and auditable.

\subsubsection{Step~2.9: Join with question metadata}
Merge the participant-question trial table with the cleaned question metadata.
At this stage, validate:
\begin{itemize}
    \item that the question type in the log matches the question type in the metadata;
    \item that the question year and field are consistent;
    \item that each presented item finds exactly one metadata row.
\end{itemize}

If any trial does not match cleanly to a question in the metadata, keep the row but flag it for manual review.

\subsubsection{Step~2.10: Derive block-order variables}
From the block labels such as \texttt{A1}, \texttt{A2}, \texttt{C1}, and so on, derive:
\begin{itemize}
    \item \texttt{block\_type}
    \item \texttt{block\_index\_within\_type}
    \item \texttt{global\_block\_order}
    \item \texttt{order\_condition}
\end{itemize}

This is important because the experimental script presents five blocks of one type first and then five blocks of the other type.
Even in a basic descriptive workflow, order should be retained as a documented variable.

\subsubsection{Step~2.11: Build the clean trial table}
Save the final cleaned participant-question dataset as:
\begin{verbatim}
06_behavior_trial_clean.csv
\end{verbatim}

This should be the main Step~2 analysis table.
It should contain one row per participant-question trial, all relevant metadata, timing variables, correctness variables, and quality flags.

\subsubsection{Step~2.12: Produce question-level summaries}
For each unique question, compute:
\begin{itemize}
    \item number of participants exposed;
    \item number of valid answers;
    \item number correct;
    \item sample accuracy;
    \item ENEM correctness;
    \item difference between sample accuracy and ENEM correctness;
    \item mean and median response time;
    \item standard deviation of response time;
    \item mean normalized timing measures;
    \item answer distribution across A--E.
\end{itemize}

Save these results as:
\begin{verbatim}
07_question_summary.csv
08_question_answer_distribution.csv
\end{verbatim}

\subsubsection{Step~2.13: Compare participant sample accuracy with ENEM correctness}
For each question, define:
\begin{verbatim}
difficulty_gap = sample_accuracy - enem_correctness
\end{verbatim}

This quantity should be interpreted descriptively:
\begin{itemize}
    \item a positive value means the sample found the item easier than the original ENEM population;
    \item a negative value means the sample found the item harder than the original ENEM population.
\end{itemize}

This comparison is useful, but it must be described carefully.
The two groups are not the same population, so this is a descriptive comparison, not a formal calibration exercise.

Save the comparison table as:
\begin{verbatim}
14_exam_vs_sample_comparison.csv
\end{verbatim}

\subsubsection{Step~2.14: Produce type-level summaries}
Aggregate the clean trial table by question type:
\begin{itemize}
    \item Abstract
    \item Concrete
\end{itemize}

For each type, compute:
\begin{itemize}
    \item number of trials;
    \item number of valid answers;
    \item mean sample accuracy;
    \item mean response time;
    \item median response time;
    \item mean ENEM correctness of the presented items;
    \item mean word count;
    \item mean sentence count;
    \item mean sentence length.
\end{itemize}

Save this table as:
\begin{verbatim}
09_type_summary.csv
\end{verbatim}

\subsubsection{Step~2.15: Produce block-level summaries}
Aggregate by:
\begin{itemize}
    \item block label;
    \item block type;
    \item global block order;
    \item trial position within block.
\end{itemize}

For each grouping, compute:
\begin{itemize}
    \item number of presented questions;
    \item number of answered questions;
    \item mean response time;
    \item mean accuracy;
    \item mean normalized timing measures.
\end{itemize}

This helps identify simple fatigue or order effects without introducing complex modeling.

Save the result as:
\begin{verbatim}
10_block_summary.csv
\end{verbatim}

\subsubsection{Step~2.16: Produce field-level and year-level summaries}
Aggregate by:
\begin{itemize}
    \item \texttt{question\_field}
    \item \texttt{question\_year}
\end{itemize}

For each group, compute:
\begin{itemize}
    \item number of questions;
    \item number of valid trials;
    \item mean accuracy;
    \item mean ENEM correctness;
    \item mean response time.
\end{itemize}

Save these tables as:
\begin{verbatim}
11_field_summary.csv
12_year_summary.csv
\end{verbatim}

\subsubsection{Step~2.17: Produce participant-level summaries}
For each participant, compute:
\begin{itemize}
    \item number of questions shown;
    \item number of valid answers;
    \item overall accuracy;
    \item overall mean response time;
    \item abstract-question accuracy;
    \item concrete-question accuracy;
    \item abstract-question mean response time;
    \item concrete-question mean response time;
    \item order condition.
\end{itemize}

This participant-level table is especially useful for later merging with fNIRS summaries.

Save it as:
\begin{verbatim}
13_participant_behavior_summary.csv
\end{verbatim}

\subsubsection{Step~2.18: Produce missingness and quality tables}
Create a missingness table that reports, for each variable and each response-status category:
\begin{itemize}
    \item count;
    \item percentage.
\end{itemize}

Save it as:
\begin{verbatim}
15_behavior_missingness.csv
\end{verbatim}

\subsection{Minimal Descriptive Analyses to Include}

\subsubsection*{A. Question-level analysis}
For each question:
\begin{itemize}
    \item sample size;
    \item sample accuracy;
    \item ENEM correctness;
    \item difficulty gap;
    \item answer distribution;
    \item mean and median response time;
    \item question length variables.
\end{itemize}

\subsubsection*{B. Abstract vs concrete analysis}
Compare abstract and concrete questions descriptively on:
\begin{itemize}
    \item sample accuracy;
    \item response time;
    \item normalized timing;
    \item average word count;
    \item average sentence count;
    \item average sentence length;
    \item average ENEM correctness.
\end{itemize}

\subsubsection*{C. Item-length analysis}
Describe how item length relates to behavior using:
\begin{itemize}
    \item scatter plot of word count versus mean response time;
    \item scatter plot of word count versus sample accuracy;
    \item optional tables sorted by longest questions and shortest questions.
\end{itemize}

\subsubsection*{D. Position analysis}
Describe how behavior varies by:
\begin{itemize}
    \item block order;
    \item question position within block;
    \item order condition.
\end{itemize}

\subsection{Plots}
The following plots are recommended for the basic Step~2 report.

\subsubsection*{Core plots}
\begin{itemize}
    \item bar plot of sample accuracy by question type;
    \item boxplot or violin-free boxplot of response time by question type;
    \item scatter plot of ENEM correctness versus sample accuracy;
    \item scatter plot of word count versus mean response time;
    \item scatter plot of word count versus sample accuracy.
\end{itemize}

\subsubsection*{Question-level plots}
\begin{itemize}
    \item bar plot of sample accuracy by question;
    \item paired comparison plot of sample accuracy and ENEM correctness by question;
    \item answer-distribution plot for each question.
\end{itemize}

\subsubsection*{Block and participant plots}
\begin{itemize}
    \item accuracy by global block order;
    \item response time by global block order;
    \item participant-level abstract versus concrete summary plot.
\end{itemize}

\subsection{Per-Question \LaTeX{} Reporting}
One modular \LaTeX{} file should be created for each question:
\begin{verbatim}
reports/step2/questions/question_<question_id>.tex
\end{verbatim}

Each file should contain:
\begin{itemize}
    \item question identifier;
    \item question metadata table;
    \item sample exposure count;
    \item answer distribution;
    \item sample accuracy;
    \item ENEM correctness;
    \item difficulty gap;
    \item response-time summary;
    \item corresponding plot or plots;
    \item notes on any data-quality issue.
\end{itemize}

\subsection{Master \LaTeX{} Report}
Create:
\begin{verbatim}
reports/step2/step2_behavior_report.tex
\end{verbatim}

This report should contain:
\begin{itemize}
    \item a short introduction to the behavioral analysis;
    \item a description of the raw files and reconstruction logic;
    \item a section on item metadata;
    \item a section on overall task performance;
    \item a section on abstract versus concrete comparisons;
    \item a section on question-level summaries;
    \item a section on item-length summaries;
    \item a section on data-quality and missingness;
    \item \texttt{\textbackslash input\{...\}} statements for modular question or group files.
\end{itemize}

\subsection{Template for One Question Report}
Each question-specific \LaTeX{} file can follow a structure such as:

\begin{verbatim}
\subsection{Question <question_id>}
\label{sec:q_<question_id>}

\paragraph{Metadata.}
Year, field, type, correct answer, ENEM correctness, word count, sentence count, sentence length.

\paragraph{Exposure and validity.}
N exposed, N answered, N valid.

\paragraph{Performance.}
Sample accuracy, ENEM correctness, difficulty gap.

\paragraph{Timing.}
Mean, SD, median for total response time and normalized timing.

\paragraph{Answer distribution.}
Counts and percentages for A, B, C, D, and E.

\paragraph{Notes.}
Any issue with missing events, inconsistent metadata, or unusual answer patterns.
\end{verbatim}

\subsection{Quality-Control Checks}
Before Step~2 is considered complete, the following checks should be run:
\begin{enumerate}
    \item every participant log is readable;
    \item each question trial has a valid participant ID and block label;
    \item event phases are reconstructed correctly;
    \item the answer event is identified correctly;
    \item chosen answers are limited to A--E;
    \item metadata merges uniquely to each question;
    \item question type in the log matches question type in the metadata;
    \item ENEM correctness values are within the interval $[0,1]$;
    \item response-time variables are non-negative;
    \item normalized timing variables are computed only when denominators are valid;
    \item denominators are reported in all summary tables.
\end{enumerate}

\subsection{What We Intentionally Do Not Do Yet}
To keep Step~2 simple and robust, the following are intentionally postponed:
\begin{itemize}
    \item no inferential tests between abstract and concrete items;
    \item no mixed-effects modeling;
    \item no item-response modeling;
    \item no formal equating between the sample and the ENEM population;
    \item no integration with fNIRS signals yet.
\end{itemize}

\subsection{Minimal Completion Criteria for Step~2}
Step~2 should be considered complete when:
\begin{itemize}
    \item all participant PsychoPy logs have been read;
    \item a clean participant-question table exists;
    \item question metadata from \texttt{filtered\_questions.json} have been merged;
    \item correctness and timing variables have been derived;
    \item item-level and type-level summary tables exist;
    \item sample accuracy has been descriptively compared with ENEM correctness;
    \item basic plots have been generated;
    \item modular \LaTeX{} outputs have been generated;
    \item all data-quality issues have been documented.
\end{itemize}

\subsection{Short Practical Summary}
In simple terms, Step~2 should transform the PsychoPy task logs into one clean behavioral dataset with one row per participant-question response, enriched with item metadata from \texttt{filtered\_questions.json}, and produce a basic descriptive analysis of:
\begin{itemize}
    \item time;
    \item correctness;
    \item abstract versus concrete differences;
    \item question difficulty;
    \item comparison with original ENEM correctness;
    \item simple item-length effects.
\end{itemize}
